% ============================================================
%  MATRICE RACI — Phase 1 : Démarrage
%  Time'Eats — Cellule PMO
% ============================================================
\documentclass[11pt,a4paper,oneside]{report}
\usepackage{../style/consulting}
\usepackage{pdflscape}
\hypersetup{pdftitle={Matrice RACI — Time'Eats PMO}}
\pagestyle{fancy}\fancyhf{}
\renewcommand{\headrulewidth}{0pt}
\fancyhead[L]{\begin{tikzpicture}[remember picture,overlay]
  \draw[cnNavy,line width=0.5pt](0,0)--(\textwidth,0);
\end{tikzpicture}\vspace{2pt}\timeeatslogo\hspace{4pt}\small\textcolor{cnSlate}{Time'Eats \textbf{·} Matrice RACI}}
\fancyhead[R]{\small\textcolor{cnSlate}{\thepage}}
\fancyfoot[C]{\tiny\textcolor{cnSlate}{Document confidentiel — Cellule PMO — CESI MAALSI 2026}}
\begin{document}\small

\begin{center}
{\LARGE\bfseries\color{cnNavy} Matrice RACI}\\[4pt]
\textcolor{cnOrange}{\rule{10cm}{1pt}}\\[4pt]
{\large\color{cnSlate} Phase 1 — Démarrage \quad|\quad Projet : \textit{[Nom du projet]}}
\end{center}

\vspace{6pt}

\begin{finding}
\textbf{Légende RACI :}
\textbf{R} = Responsable (réalise) \quad
\textbf{A} = Accountable (valide, rend compte) \quad
\textbf{C} = Consulté (sollicité avant décision) \quad
\textbf{I} = Informé (tenu au courant)\\
\textit{Règle : une seule personne \textbf{A} par activité. Au moins un \textbf{R} par activité.}
\end{finding}

\vspace{6pt}

\begin{landscape}
\renewcommand{\arraystretch}{1.4}
\begin{tabularx}{\linewidth}{L{5.5cm} C{2cm} C{2cm} C{2cm} C{2cm} C{2cm} C{2cm} C{2cm}}
\toprule
\rowcolor{cnNavy}
\th{Activité / Livrable} & \th{Chef de projet} & \th{DSI} & \th{Direction Générale} & \th{PMO} & \th{ESN / Presta} & \th{MOA / Métier} & \th{DPO / RH} \\
\midrule
\multicolumn{8}{l}{\cellcolor{cnNavy!15}\textbf{\textcolor{cnNavy}{Phase 1 — Démarrage}}} \\
\midrule
Rédiger la note de cadrage        & R & A & C & I & -- & C & -- \\
\rowcolor{cnIce}
Rédiger la charte projet          & R & A & I & I & -- & C & -- \\
Constituer l'équipe projet        & R & A & I & C & -- & -- & C \\
\rowcolor{cnIce}
Organiser le kick-off             & R & A & I & I & C & C & -- \\
\midrule
\multicolumn{8}{l}{\cellcolor{cnNavy!15}\textbf{\textcolor{cnNavy}{Phase 2 — Planification}}} \\
\midrule
Rédiger le PMP                    & R & A & I & C & C & C & -- \\
\rowcolor{cnIce}
Construire le WBS                 & R & A & -- & C & C & C & -- \\
Élaborer le planning Gantt        & R & A & -- & C & C & -- & -- \\
\rowcolor{cnIce}
Chiffrer le budget                & R & A & A & C & C & -- & -- \\
Identifier et évaluer les risques & R & C & -- & A & C & C & -- \\
\rowcolor{cnIce}
Rédiger le plan de communication  & R & A & I & C & -- & C & -- \\
\midrule
\multicolumn{8}{l}{\cellcolor{cnNavy!15}\textbf{\textcolor{cnNavy}{Phase 3 — Exécution}}} \\
\midrule
Piloter les prestataires ESN      & R & A & -- & C & R & -- & -- \\
\rowcolor{cnIce}
Rédiger les CR hebdomadaires      & R & I & -- & I & I & I & -- \\
Gérer les demandes de changement  & R & A & I & C & C & C & -- \\
\rowcolor{cnIce}
Suivre les anomalies recette      & R & C & -- & I & R & C & -- \\
\midrule
\multicolumn{8}{l}{\cellcolor{cnNavy!15}\textbf{\textcolor{cnNavy}{Phase 4 — Maîtrise}}} \\
\midrule
Produire le tableau de bord       & R & A & I & C & -- & I & -- \\
\rowcolor{cnIce}
Animer le COPIL mensuel           & R & A & I & C & I & I & -- \\
Émettre un rapport d'écarts       & R & A & I & C & C & -- & -- \\
\midrule
\multicolumn{8}{l}{\cellcolor{cnNavy!15}\textbf{\textcolor{cnNavy}{Phase 5 — Clôture}}} \\
\midrule
Organiser la recette fonctionnelle & R & C & -- & I & R & A & -- \\
\rowcolor{cnIce}
Signer le PV de transfert         & R & A & I & I & I & -- & -- \\
Rédiger la fiche REX              & R & A & -- & A & C & C & -- \\
\bottomrule
\end{tabularx}
\end{landscape}

\end{document}
